\section{Conclusion}
\label{sec:conclusion}

We have introduced \ams{} (Adaptive Multi-scale MCMC), a Search-layer
algorithm that wraps \ms{} (G.11) with a Robbins-Monro adaptation loop that
automatically tunes the coarse-move probability $\alpha$.
The algorithm is simple: every $\Delta = 50$ steps, observe the recent coarse
acceptance rate, compare to the target $\tau = 0.30$, and update $\alpha$ by
$\gamma_t = \gamma_0/t^{3/4}$.
No other aspect of the chain changes.

\paragraph{Key findings.}
\begin{itemize}
  \item \textbf{State-specific convergence.}
    \ams{} converges to distinct state-specific equilibria: $\alpha^* \approx
    0.33$ for Texas (highly county-structured), $0.28$ for North Carolina
    (well-matched to the default), and $0.17$ for Iowa (few districts, small
    ratio).
    These differ from the one-size-fits-all default of 0.30.

  \item \textbf{No regression from adaptation.}
    Autocorrelation of $\EC$ at lag 100 is not worse than fixed-$\alpha$
    \ms{} using the oracle-optimal alpha in any state tested.
    For Iowa and California---states where 0.30 is a poor choice---\ams{}
    provides a clear improvement over the fixed default.

  \item \textbf{Near-oracle performance.}
    ACF from \ams{} is within 0.01--0.02 of the oracle-optimal fixed-$\alpha$
    chain (Table~\ref{tab:oracle}), achieved without any manual tuning or
    pilot runs.

  \item \textbf{Negligible overhead.}
    The adaptation loop is $O(1)$ per step in amortised cost; runtime is
    within 2\% of fixed-$\alpha$ \ms{} for all states.
\end{itemize}

\paragraph{Position in the B-series landscape.}
\ams{} is a drop-in replacement for \texttt{--search multiscale} at the
Search (SeedCompositor) layer.
It is orthogonal to all Structure-layer choices (\geosec{}, \areasec{},
prime-factor, \sa{}) and all WeightsOverride choices.
For practitioners, the recommendation is straightforward:
\begin{itemize}
  \item Use \texttt{--search multiscale-adaptive} whenever you would have
    used \texttt{--search multiscale}.
    The adaptation cost is negligible and the calibration benefit is real.
  \item Use \texttt{--search multiscale --multiscale-alpha X} only when you
    have strong prior knowledge that $X$ is the correct $\alpha$ for this
    specific state, or for reproducibility of a specific historical run.
\end{itemize}

\paragraph{Distributional footprint and partisan seat stability.}
The distributional footprint of \ams{} --- the space of plans it visits ---
depends on the converged $\alpha^*$, which is state-specific.
States where \ams{} converges to higher $\alpha$ (more coarse moves) visit
a broader spatial range of plans; this does not introduce partisan bias since
coarse moves are geographically symmetric, but practitioners should run the
full $p$-sweep (as recommended for all ensemble methods) to characterise
seat outcome sensitivity.

\paragraph{Audit and reproducibility.}
The alpha trace, $\gamma_0$, initial alpha, and final alpha are all recorded
in the run manifest.
A verifier can independently reproduce the alpha trace from the base seed
and step seeds using the MSC\_STEP\_ prefix formula.
The adaptation introduces no new randomness: it is a deterministic function
of the acceptance events, which are themselves determined by the step seeds.

\paragraph{Open questions.}
Several extensions are deferred to future work:

\begin{enumerate}
  \item \textbf{Frozen vs.\ continued adaptation.}
    The current implementation adapts throughout the run.
    An alternative is to freeze $\alpha$ once convergence is detected
    (e.g., when $|\alpha_t - \alpha_{t-1}| < \epsilon$ for several consecutive
    rounds), giving a cleaner theoretical guarantee (the frozen chain
    satisfies detailed balance exactly).
    Continued adaptation is more responsive to non-stationarity; frozen
    alpha is more reproducible.

  \item \textbf{State-dependent target $\tau$.}
    The current spec uses a fixed $\tau = 0.30$ for all states.
    One could set $\tau$ proportional to the county/tract ratio, so that
    states with higher ratios target higher coarse-move frequencies.
    This would require a prior model for $\tau(\text{ratio})$ that does not
    yet exist.

  \item \textbf{Alpha trace compression.}
    For $T = 10{,}000$ and $\Delta = 50$, the alpha trace has 200 entries---
    modest but potentially significant for large ensemble runs.
    Run-length encoding of near-converged traces would reduce storage while
    preserving full verifiability.

  \item \textbf{Extension to multi-chain ensembles.}
    When \ams{} is composed with \cs{} or \be{}, each chain runs its own
    adaptation independently.
    A shared adaptation pooling coarse acceptance rates across chains might
    converge faster; the interaction with the SeedCompositor design is an
    open question.
\end{enumerate}

\paragraph{Implementation note.}
\ams{} is implemented as \texttt{AdaptiveMultiScaleChain} in the
\texttt{redist-multiscale} crate, wrapping \texttt{MultiScaleChain} by
composition.
The CLI flag \texttt{--search multiscale-adaptive} enables it; the YAML
key is \texttt{search: multiscale-adaptive}.
All other platform interfaces (\texttt{redist label-analyze},
\texttt{redist label-verify}, the manifest schema) are unchanged.
